People learn whom to trust from feedback. When advice can be checked against the outcome, people track how trustworthy each advisor is, and their choices are consistent with a Bayesian learner that updates faster when an advisor's accuracy is volatile \citep{behrens2008associative}; advisor reputations form quickly from a track record \citep{yaniv2000advice,bonaccio2006advice}. We adapt this human-neuroscience paradigm and its computational modeling approach to evaluate vision-language models (VLMs).

Agents built on large language models (LLMs) and VLMs face the same problem when they combine tool outputs, retrieved documents or messages from other agents whose reliability can change; for VLM-based agents, some of these messages arrive as images. Prior work has tested source weighting under fixed reliability \citep{ashkinaze2026information} and test-time reliability inversion \citep{zhou2026epistemic}, studied trust formation and recovery after scripted peer failures \citep{chen2026trust}, and proposed memories that track reliability outside the model \citep{feng2026sigmamem}. We ask how models learn in context, without parameter updates, from trial-by-trial feedback across repeated reversals of advisor accuracy, and how the format of that history affects their choices.

The format of past evidence matters: conversational LLMs can choose differently when prospects are stated explicitly than when they are given as histories of outcomes \citep{ge2026dhgap}, and summarized histories can improve exploration in bandits \citep{krishnamurthy2024explore}. Learning whom to trust additionally requires deriving advisor correctness by comparing conflicting reports with observed outcomes. It involves three computations: (i) judging which advisor was correct on each trial, (ii) integrating these judgments across trials, and (iii) using the estimate; in the human study, (i) yields the observation from which a prediction error about the advisor is computed \citep{behrens2008associative}. Five VLMs receive the same history in formats that state the result of (i), or of (i) and (ii), to different degrees (Fig.~\ref{fig:teaser}), with current reports as text or images. Because formats also change length and salience, we test which formats support good choices without assuming separable internal stages. Our contributions are:
\begin{itemize}
  \item[$\bullet$] An advice-taking evaluation adapted from human neuroscience, with text and image reports.
  \item[$\bullet$] For five open VLMs answering directly, \CorrectnessOnly{} recovers most of the \Summary{}$-$\Raw{} gap, and a masking test predicted before held-out evaluation shows that the visible outcome label drives repetition of the previous outcome.
  \item[$\bullet$] A boundary condition: two proprietary reasoning models have \Raw{} accuracy similar to that of a Bayesian volatility learner, so the history-format effect depends on the system.
  \item[$\bullet$] A descriptive human--model comparison in a Psych-101 changing bandit: participants' estimated forgetting rose with the hazard rate (the per-trial probability of a payoff change), whereas the open models' estimates changed little.
\end{itemize}

% Related Work merges (keep blue, drop red):
% Cognitive evaluation: "... found asymmetric or biased updating and rigid choice \citep{schubert2024incontext,hayes2024biased,wang2026rigidity}, as well as premature greedy action selection \citep{schmied2025greedy}; reversal-learning schedules have been compared across LLMs \citep{wang2026comparative,li2025reflectionbench}."
% Volatility: "..., and arbitration between social and individual learning has been modeled using the relative precision of their predictions \citep{diaconescu2020neural}."
